\section{The architecture ladder}

Holding the metric fixed, we walk the representation class $\mathcal{F}$ of Eq.~\eqref{eq:tuple} up four rungs. Rungs 1--3 are analysed in general two-axis form and evidenced synthetically in Section~4; rung 4 is evidenced on the short-time Fourier transform (STFT) band--frame plane of a public benchmark, where $\Sigma^{-1}$ is diagonal in frequency by measurement (Section~4).

\paragraph{Rung 1: rank-one fixed template (optimal filtering).}
Fixing the signal to a known template $s_0$, the likelihood-optimal reconstruction is the optimal filter (OF) amplitude of Section~2 \citep{gatti1986processing,turin1960introduction}. The representation class is a single direction, and the metric question has a closed-form answer.

\paragraph{Rung 2: unrestricted rank-$k$ linear.}
Replacing the fixed template with a learned rank-$k$ subspace gives covariance-weighted PCA (CW-PCA), which is also the tied noise-aware linear autoencoder: both minimize the likelihood residual over the full Grassmannian $\Gr(k,n_1n_2)$ of a two-axis measurement (Proposition~\ref{prop:bridge}), with reconstruction $\hat s=PP^\top\Sigma^{-1}x$.

\paragraph{Rung 3: Kronecker-separable two-axis projection (NFPA).}
NFPA restricts rung 2's projection to separable factors of a two-axis measurement $Z_i\in\R^{n_1\times n_2}$ --- sensor channel $\times$ time in the controlled synthetic study, frequency band $\times$ time frame on the public benchmark. It learns orthonormal bases $U_1\in\R^{n_1\times k_1}$ and $U_2\in\R^{n_2\times k_2}$, with code $C_i=U_1^\top Z_iU_2$ and reconstruction $\hat Z_i=U_1U_1^\top Z_iU_2U_2^\top$, so that
\begin{equation}
  \mathrm{vec}(\hat Z_i)=\bigl(U_2U_2^\top\otimes U_1U_1^\top\bigr)\mathrm{vec}(Z_i),
  \label{eq:kron-projector}
\end{equation}
and the joint representation is $\range(U_1)\otimes\range(U_2)$. The admissible family is the embedded product
\begin{equation}
  \Gr(k_1,n_1)\times\Gr(k_2,n_2)
  \hookrightarrow \Gr(k_1k_2,n_1n_2).
  \label{eq:product-grassmann}
\end{equation}

\begin{proposition}[Kronecker-restricted bridge]
\label{prop:restricted}
NFPA minimizes the CW-PCA/tied-autoencoder objective restricted to Kronecker-separable two-axis subspaces. Its optimum equals the unrestricted CW-PCA optimum if and only if at least one rank-$k_1k_2$ CW-PCA minimizer has such a separable subspace.
\end{proposition}

\paragraph{Proof sketch.}
Equation~\eqref{eq:kron-projector} is the orthogonal projector onto $\range(U_1)\otimes\range(U_2)$ in whitened coordinates, so NFPA is the whitened rank-$k_1k_2$ objective restricted to Eq.~\eqref{eq:product-grassmann}; equality holds precisely when that set contains an unrestricted minimizer, and otherwise the gap is the cost of the architecture prior, not a failure of the metric. This is a covariance-aware member of the multilinear-PCA family \citep{lu2008mpca}, not a new factorization. When either factor is full rank the restriction is inactive and rung 3 reduces to rung 2; at $k_1=k_2=1$ it recovers rung 1 (Corollary~\ref{cor:rankone}, Appendix~\ref{app:proofs}).

\paragraph{Rung 4: factored attention.}
Rung 3's two-factor restriction, instantiated on the band--frame plane, reads as $\range(U_f)\otimes\range(U_m)$ --- separable frequency-band and time-frame factors; rung 4's transformer performs intra-band temporal attention over frames, then cross-band attention. These are the same factorization at different levels of nonlinearity, in the same coordinates, and $\Sigma^{-1}$ is diagonal in exactly the basis the first factor indexes: on this benchmark a stationary-PSD model fits the readout noise ($\chi^2/\mathrm{dof}=1.0002$, all gates passed; Section~4). The rung-4 model is a transformer \citep{vaswani2017attention} over frame patches \citep{nie2023patchtst}, with rotary embeddings on the temporal axis \citep{su2024roformer} and learned embeddings across bands; the metric question is posed empirically, training identical architecture, data, and budget on two objectives differing only in the metric (Section~4). Injecting the optimal-filter statistic into the forward pass follows prior work on calorimetric event reconstruction \citep{priorThesis} and serves as a secondary contrast.
